\section*{TÓM TẮT}
\phantomsection\addcontentsline{toc}{section}{\numberline{}TÓM TẮT}
Đồ án nghiên cứu phương pháp ước lượng thời gian hiện diện của nhiều nhân viên trong khu vực làm việc từ một camera CCTV cố định. Thời gian làm việc hiệu lực được định nghĩa là hiện diện hợp lệ sau khi áp dụng chính sách điều chỉnh thời gian sử dụng điện thoại. Đại lượng này không được dùng để suy luận năng suất, mức độ tập trung hoặc chất lượng công việc.

Giải pháp đề xuất kết hợp YOLOv8 để phát hiện người và điện thoại, ByteTrack để duy trì quỹ đạo, vùng quan tâm dạng đa giác để liên kết vị trí làm việc và vector đặc trưng khuôn mặt để hỗ trợ xác minh danh tính. Hai máy trạng thái theo thời gian tạo khoảng hiện diện và phiên điện thoại. Bộ tổng hợp lấy giao với lịch, loại khoảng không quan sát được và áp dụng hạn mức theo ca. Kiến trúc quản trị sử dụng FastAPI, PostgreSQL, Redis và WebSocket, kèm phân quyền, phiên bản chính sách và kho sự kiện phục vụ truy vết.

Phần thực nghiệm tính toán được thực hiện trên bộ tổng hợp tham chiếu với đầu vào tổng hợp. Trong 24 kịch bản xác định, 20 trường hợp hợp lệ khớp đáp án và bốn trường hợp sai bị từ chối. Bộ 2.000 đầu vào ngẫu nhiên được đối chiếu với một cách biểu diễn độc lập theo từng giây, không ghi nhận chênh lệch thời lượng. Các phép biến đổi thứ tự, nhân đôi, dịch mốc và chia khoảng giữ nguyên kết quả. Bộ số đếm, 12 phiên điện thoại và 30 ca có nhiễu được sử dụng riêng để kiểm tra công thức chỉ số và phân tích độ nhạy theo giả định.

Kết quả hỗ trợ tính đúng đắn của bộ tính trong phạm vi kiểm thử đã xét. Dữ liệu chưa được hiệu chỉnh theo video thực địa, vì vậy chưa dùng để xác nhận độ chính xác nhận biết hoặc hiệu năng của hệ thống CCTV. Báo cáo trình bày quy trình đánh giá trên video có nhãn tham chiếu và các kiểm thử tích hợp cần thực hiện tiếp theo.

\textbf{Từ khóa:} thị giác máy tính, YOLOv8, ByteTrack, vùng quan tâm, thời gian hiện diện, sử dụng điện thoại, đại số khoảng, kiểm thử phần mềm.
\clearpage
